% Chapter 4: MLOps Toolchain, Databases & Multi-Tier Storage Architecture

\begin{frame}[plain]
    \vfill
    \centering
    {\footnotesize \textbf{\textcolor{deepblue}{\textsc{TaxonVision-MLOps Architecture}}}} \\[0.35cm]
    {\large \textbf{\textcolor{signalorange}{CHAPTER 04}}} \\[0.25cm]
    {\huge \textbf{MLOps Toolchain \& Storage Architecture}} \\[0.35cm]
    {\normalsize \textcolor{black!70}{Polyglot Persistence, Tensor Partitions \& Object Storage Interlock}} \\[0.5cm]
    \textcolor{deepblue}{\rule{0.55\linewidth}{1.2pt}} \\[0.5cm]
    \statusdone
    \vfill
    \script{In Chapter 4, we examine our multi-tier storage topology: how we interlock Apache Arrow columnar formats, PyTorch tensor caching, SQLite state databases, and a three-tier S3 object hierarchy to achieve microsecond training throughput.}
    \qa{What is the core storage philosophy of TaxonVision?}{Polyglot persistence: selecting specialized data stores for each access pattern rather than forcing an end-to-end relational database or raw filesystem upon heterogeneous ML workloads.}
\end{frame}

% ---------------------------------------------------------------------
% Slide 2: Technology & Storage Matrix (Part 1: Ingestion & Features)
% ---------------------------------------------------------------------
\begin{frame}{Storage Technology Matrix}{Tabular Manifests, Tensor Caching \& Optimization State}
    \begin{center}\scriptsize
        \begin{tabular}{p{2.4cm} p{1.8cm} p{2.8cm} p{4.7cm}}
            \toprule
            \textbf{Technology} & \textbf{Storage Tier} & \textbf{Location / URI} & \textbf{Selection Rationale vs Alternatives} \\
            \midrule
            \textbf{Apache Parquet / Polars} & Host NVMe SSD & \texttt{data/manifests/} \newline \texttt{*.parquet} & Columnar zero-copy IPC and $10\times$ compression over CSV/JSON; lazy streaming queries without RAM exhaustion. \\
            \midrule
            \textbf{PyTorch Split Tensors (\texttt{.pt})} & Host NVMe \& DVC S3 & \texttt{data/embeddings/} \newline \texttt{*.pt} & Pre-computed foundation embeddings eliminate redundant JPEG decoding, accelerating training by $50\times$. \\
            \midrule
            \textbf{SQLite on PVC} & k3s Persistent Volume & \texttt{sqlite:////workspace/} \newline \texttt{optuna/optuna\_studies.db} & Zero-maintenance ACID transactions for Bayesian study states without managing an external Postgres cluster. \\
            \bottomrule
        \end{tabular}
    \end{center}

    \vspace{0.05cm}
    \statusdone

    \script{This matrix outlines our data tier. By pairing Apache Parquet for metadata with pre-computed PyTorch tensor splits, our training loops avoid the massive disk I/O bottleneck of decoding millions of JPEG images on every epoch.}
    \qa{Why not store metadata directly in SQLite?}{Parquet provides column-oriented vector compression and direct Arrow interop with Polars, achieving gigabyte-per-second scan speeds that row-based SQLite cannot match for analytical batches.}
\end{frame}

% ---------------------------------------------------------------------
% Slide 3: Technology & Storage Matrix (Part 2: Remotes & Telemetry)
% ---------------------------------------------------------------------
\begin{frame}{Storage Technology Matrix}{Object Storage, Experiment Lineage \& Observability TSDB}
    \begin{center}\scriptsize
        \begin{tabular}{p{2.4cm} p{1.8cm} p{2.8cm} p{4.7cm}}
            \toprule
            \textbf{Technology} & \textbf{Storage Tier} & \textbf{Location / URI} & \textbf{Selection Rationale vs Alternatives} \\
            \midrule
            \textbf{SeaweedFS S3} & In-Cluster Edge Cache & \texttt{http://seaweedfs-} \newline \texttt{s3:8333} & Ultra-low latency S3 API with O(1) disk lookup; $4\times$ less memory overhead than MinIO or Ceph on bare metal. \\
            \midrule
            \textbf{DagsHub DVC Remote S3} & Cloud Immutable Remote & \texttt{s3://dvc/} \newline \textsf{via DagsHub} & Free collaborative DVC remote storage bound directly to Git commit lineage without managing cloud buckets. \\
            \midrule
            \textbf{MLflow Registry} & Cloud Tracking Server & \texttt{https://dagshub.com/} \newline \texttt{.../mlflow} & Full model lifecycle governance, explicit tensor schemas, and automated promotion aliases (\texttt{@Production}, \texttt{@Challenger}). \\
            \midrule
            \textbf{Prometheus TSDB} & In-Memory / Host Disk & Scrapes \texttt{api:8000/metrics} & Cloud-Native CNCF time-series engine tracking p95 latency, error rates, and Wasserstein feature drift. \\
            \bottomrule
        \end{tabular}
    \end{center}

    \vspace{0.05cm}
    \statusdone

    \script{Here we review the remote storage and telemetry tiers: in-cluster SeaweedFS for microsecond DVC caching, DagsHub for collaborative release immutability, and Prometheus for real-time drift telemetry.}
    \qa{Why choose SeaweedFS over MinIO for in-cluster storage?}{SeaweedFS separates volume storage from directory lookup, achieving constant-time O(1) disk reads with significantly lower CPU and RAM footprints on resource-constrained bare-metal nodes.}
\end{frame}

% ---------------------------------------------------------------------
% Slide 4: Tabular & Analytical Engine Interlock
% ---------------------------------------------------------------------
\begin{frame}{Analytical Storage}{The Tabular Interlock: Arrow Columnar Memory, Parquet \& Polars}
    \begin{columns}[T]
        \begin{column}{0.48\textwidth}
            \begin{block}{Apache Arrow \& Parquet Storage}
                \begin{itemize}\footnotesize
                    \item \textbf{Columnar Layout:} Contiguous in-memory memory buffers allow SIMD vectorized filtering across millions of observation rows.
                    \item \textbf{Dictionary Encoding:} Categorical attributes (e.g. \texttt{license}, \texttt{taxon\_id}, \texttt{quality\_grade}) compress to 1-byte indices.
                    \item \textbf{Predicate Pushdown:} Parquet page statistics skip unread row groups during license auditing without reading raw bytes into memory.
                \end{itemize}
            \end{block}
        \end{column}
        \begin{column}{0.48\textwidth}
            \begin{block}{Polars Multi-Threaded Engine}
                \begin{itemize}\footnotesize
                    \item \textbf{Rust Core:} Executes analytical queries in parallel across all available CPU cores without Python GIL locking.
                    \item \textbf{Lazy Evaluation (\texttt{LazyFrame}):} Optimizes query filter graphs before execution, merging projections and predicates.
                    \item \textbf{DuckDB Interop:} Enables zero-copy in-process SQL queries directly against Parquet manifests without running a database server.
                \end{itemize}
            \end{block}
        \end{column}
    \end{columns}

    \vspace{0.1cm}
    \statusdone

    \script{Our tabular data engine leverages Apache Arrow, Parquet, and Polars. By operating entirely in-memory with columnar zero-copy buffers, we audit licensing and stratify manifests across thousands of observations in milliseconds.}
    \qa{How does predicate pushdown benefit the license audit pipeline?}{It inspects Parquet column chunk headers to skip non-compliant or non-research-grade records at the storage level, reducing disk I/O by over 80 percent.}
\end{frame}

% ---------------------------------------------------------------------
% Slide 5: Epistemic Diamond: 4-Way Tensor Partitioning
% ---------------------------------------------------------------------
\begin{frame}{Tensor Caching}{The Epistemic Diamond: 4-Way Zero-Leakage Tensor Partitioning}
    \begin{columns}[T]
        \begin{column}{0.48\textwidth}
            \begin{block}{The Four Strategic Tensors (\texttt{.pt})}
                \begin{itemize}\footnotesize
                    \item \textbf{\texttt{train\_embeddings.pt}} ($60\%$): Feeds the linear classification head during Class-Balanced focal loss optimization.
                    \item \textbf{\texttt{val\_embeddings.pt}} ($15\%$): Guides Optuna Bayesian hyperparameter optimization and early stopping callbacks.
                    \item \textbf{\texttt{cal\_embeddings.pt}} ($15\%$): Exclusively reserved for Split Conformal Prediction calibration ($\hat{q}$).
                    \item \textbf{\texttt{test\_embeddings.pt}} ($10\%$): Pure holdout set for final Pareto latency/accuracy benchmarking.
                \end{itemize}
            \end{block}
        \end{column}
        \begin{column}{0.48\textwidth}
            \begin{block}{Mathematical Safety Invariant}
                \begin{itemize}\footnotesize
                    \item \textbf{Strict Isolation:} Calibration tensors must never overlap with validation or training splits.
                    \item \textbf{Exchangeability Guarantee:} The calibration split satisfies the i.i.d. exchangeability property, guaranteeing that:
                        $$\mathbb{P}\left(Y_{\text{test}} \in C(X_{\text{test}})\right) \ge 1 - \alpha$$
                    \item \textbf{Zero Retraining Overhead:} Training heads takes $< 2\text{ seconds}$ on pre-computed embeddings.
                \end{itemize}
            \end{block}
        \end{column}
    \end{columns}

    \vspace{0.1cm}
    \statusdone

    \script{This slide details our 4-way stratified embedding split: the Epistemic Diamond. By isolating independent calibration and holdout test partitions before training, we guarantee the mathematical validity of our conformal coverage bounds.}
    \qa{Why must cal\_embeddings.pt be strictly isolated from val\_embeddings.pt?}{If the calibration set were used for model selection or tuning, the non-conformity scores would be biased downwards, violating exchangeability and breaking the 95\% coverage guarantee.}
\end{frame}

% ---------------------------------------------------------------------
% Slide 6: Hyperparameter Optimization Database (Optuna SQLite)
% ---------------------------------------------------------------------
\begin{frame}{Optimization Database}{Optuna Bayesian Study Persistence on Kubernetes PVC}
    \begin{columns}[T]
        \begin{column}{0.48\textwidth}
            \begin{block}{SQLite Persistence Architecture}
                \begin{itemize}\footnotesize
                    \item \textbf{Physical Storage URI:} \texttt{sqlite:////workspace/optuna/} \texttt{optuna\_studies.db}.
                    \item \textbf{Kubernetes PVC:} Bound to \texttt{mlops-shared-pvc} backed by hostPath storage on \texttt{hive-mind} NVMe SSD.
                    \item \textbf{Multi-Trial Resiliency:} Studies persist across ephemeral Jenkins pod terminations; crashed runs resume seamlessly.
                    \item \textbf{Search Algorithms:} Tree-structured Parzen Estimator (TPE) with ASHA median early stopping.
                \end{itemize}
            \end{block}
        \end{column}
        \begin{column}{0.48\textwidth}
            \begin{block}{Interactive Study Inspection}
                \begin{itemize}\footnotesize
                    \item \textbf{Local Dashboard:}
                        \texttt{just optuna-local} \\
                        \scriptsize Starts local \texttt{optuna-dashboard} on port 8081.
                    \item \textbf{Cluster Dashboard:}
                        \texttt{just optuna-k8s} \\
                        \scriptsize Port-forwards live cluster service to localhost:8081.
                    \item \textbf{Diagnostics:} Visualizes parameter importances, Pareto frontiers, and trial pruning timelines in real time.
                \end{itemize}
            \end{block}
        \end{column}
    \end{columns}

    \vspace{0.1cm}
    \statusdone

    \script{Optuna study states are backed by an ACID SQLite database on a persistent Kubernetes volume. This allows long-running Bayesian sweeps to survive pod restarts and enables live monitoring via optuna-dashboard.}
    \qa{How do developer workstations connect to cluster Optuna studies?}{The Justfile provides the optuna-k8s recipe, which establishes a secure kubectl port-forward directly to the cluster Optuna dashboard pod on port 8081.}
\end{frame}

% ---------------------------------------------------------------------
% Slide 7: Three-Tier S3 Object Storage Interlock
% ---------------------------------------------------------------------
\begin{frame}{Object Storage Architecture}{The Three-Tier S3 Interlock: Cloud Open Data to In-Cluster Edge}
    \begin{columns}[T]
        \begin{column}{0.32\textwidth}
            \begin{block}{Tier 1: Cloud Open Data}
                \begin{itemize}\scriptsize
                    \item \textbf{Endpoint:} AWS Open Data Registry (\texttt{s3://inaturalist-open-data/})
                    \item \textbf{Payload:} Millions of raw citizen-science specimen photographs.
                    \item \textbf{Access Pattern:} Public read-only HTTP streaming; zero egress fees.
                \end{itemize}
            \end{block}
        \end{column}
        \begin{column}{0.32\textwidth}
            \begin{block}{Tier 2: In-Cluster Edge}
                \begin{itemize}\scriptsize
                    \item \textbf{Endpoint:} SeaweedFS S3 (\texttt{http://seaweedfs-s3:8333})
                    \item \textbf{Payload:} Stratified embedding partitions \& local DVC cache.
                    \item \textbf{Access Pattern:} Microsecond reads on \texttt{hive-mind} NVMe SSD.
                \end{itemize}
            \end{block}
        \end{column}
        \begin{column}{0.32\textwidth}
            \begin{block}{Tier 3: Release Remote}
                \begin{itemize}\scriptsize
                    \item \textbf{Endpoint:} DagsHub S3 Remote (\texttt{s3://dvc/})
                    \item \textbf{Payload:} Audited release manifests, head weights, INT8 models.
                    \item \textbf{Access Pattern:} CI push/pull tied to Git commits; collaborative access.
                \end{itemize}
            \end{block}
        \end{column}
    \end{columns}

    \vspace{0.1cm}
    \begin{block}{Bandwidth \& Cost Optimization}
        Raw photographs stream once from AWS Open Data; extracted 768-d feature vectors are cached forever in Tier 2/3, reducing external bandwidth consumption by $99.8\%$.
    \end{block}

    \script{Our object storage architecture spans three tiers: streaming raw images from AWS Open Data, caching extracted tensors in bare-metal SeaweedFS, and committing final releases to DagsHub S3.}
    \qa{How does this architecture eliminate S3 egress costs?}{Raw image downloads occur over AWS public datasets without authentication charges, while repetitive pipeline iterations execute against in-cluster SeaweedFS on local NVMe storage.}
\end{frame}

% ---------------------------------------------------------------------
% Slide 8: Experiment Tracking & Model Registry (MLflow)
% ---------------------------------------------------------------------
\begin{frame}{Model Governance}{Experiment Tracking \& Model Registry with Explicit Schemas}
    \begin{columns}[T]
        \begin{column}{0.48\textwidth}
            \begin{block}{Cryptographic Lineage Binding}
                \begin{itemize}\footnotesize
                    \item \textbf{Context Scope:} Managed via \texttt{mlflow\_run\_scope()} in \texttt{tracking.py}.
                    \item \textbf{Metadata Binding:} Automatically records Git commit SHA, DVC dataset hash, and full parameter payloads.
                    \item \textbf{Metric Logging:} Logs training loss, validation accuracy, Macro F1, and Macro PR-AUC per epoch.
                \end{itemize}
            \end{block}
        \end{column}
        \begin{column}{0.48\textwidth}
            \begin{block}{Explicit Model Signatures}
                \begin{itemize}\footnotesize
                    \item \textbf{Input Schema:}
                        \texttt{TensorSpec(float32, (-1, D), name="backbone\_embeddings")}
                    \item \textbf{Output Schema:}
                        \texttt{TensorSpec(float32, (-1, K), name="taxonomic\_logits")}
                    \item \textbf{Registry Aliases:}
                        \begin{itemize}\scriptsize
                            \item \texttt{@Production}: Current active serving model.
                            \item \texttt{@Challenger}: Meets/beats baseline PR-AUC.
                            \item \texttt{Archived}: Sub-par candidates.
                        \end{itemize}
                \end{itemize}
            \end{block}
        \end{column}
    \end{columns}

    \vspace{0.1cm}
    \statusdone

    \script{MLflow binds experiment tracking to our model registry. We construct explicit ModelSignatures using MLflow TensorSpec schemas, preventing model inference crashes from shape mismatches.}
    \qa{How does MLflow enforce automated model promotion?}{When a model is registered, the tracker compares its validation PR-AUC against the production baseline. If superior, it tags the candidate as Challenger; otherwise it archives it.}
\end{frame}

% ---------------------------------------------------------------------
% Slide 9: Hermetic Environment & Quality Toolchain (Pixi)
% ---------------------------------------------------------------------
\begin{frame}{Environment Toolchain}{Pure Pixi Architecture: Zero Virtual Environments \& Rattler Cache}
    \begin{columns}[T]
        \begin{column}{0.48\textwidth}
            \begin{block}{The Anti-Pattern: Virtual Environments}
                \begin{itemize}\footnotesize
                    \item \textbf{System Pollution:} Dynamic \texttt{.venv/} directories diverge across machines, leading to ``works on my machine'' syndrome.
                    \item \textbf{Binary Incompatibility:} Raw pip installs struggle with compiled CUDA binaries, cuDNN, and ONNX Runtime backends.
                    \item \textbf{Rule 00 Mandate:} The codebase strictly bans \texttt{pip}, \texttt{poetry}, \texttt{uv}, and the creation of \texttt{.venv/}.
                \end{itemize}
            \end{block}
        \end{column}
        \begin{column}{0.48\textwidth}
            \begin{block}{The Pixi Solution}
                \begin{itemize}\footnotesize
                    \item \textbf{Hermetic Lockfile (\texttt{pixi.lock}):} Pins all native binary libraries, shared objects, and Python packages cryptographically.
                    \item \textbf{Host Rattler Cache:} Mounts a 14 GB NVMe cache (\texttt{/root/.cache/rattler/cache}) into Kubernetes agents, provisioning environments in $< 5\text{ seconds}$.
                    \item \textbf{Target Environments:} Lean CPU (\texttt{ci-dev}) for CI vs CUDA 12 GPU acceleration (\texttt{dev}) for training.
                \end{itemize}
            \end{block}
        \end{column}
    \end{columns}

    \vspace{0.1cm}
    \statusdone

    \script{We completely reject traditional virtual environments. Using Pixi and Rattler, our dependencies are resolved once into a cryptographic lockfile and cached globally across bare-metal NVMe storage.}
    \qa{How does Pixi speed up Kubernetes ephemeral agent startup?}{By mounting the host Rattler cache directory directly into runner pods, Pixi hardlinks pre-downloaded binary packages in seconds without downloading anything from the internet.}
\end{frame}

% ---------------------------------------------------------------------
% Slide 10: Strict Quality Assurance Quartet
% ---------------------------------------------------------------------
\begin{frame}{Quality Assurance}{The Static Verification Quartet: Ruff, Mypy, Complexipy \& Pytest}
    \begin{columns}[T]
        \begin{column}{0.48\textwidth}
            \begin{block}{Static Analysis \& Complexity}
                \begin{itemize}\footnotesize
                    \item \textbf{Ruff (Linter \& Formatter):} Executes 500+ linting rules in milliseconds; formats code with strict flake8/black compliance.
                    \item \textbf{Mypy (Strict Static Typing):} Verifies type safety across all 66 modules with zero dynamic reflection (\texttt{hasattr}/\texttt{getattr} banned).
                    \item \textbf{Complexipy (Cognitive Complexity):} Guarantees every function maintains a complexity score $< 15$, preventing spaghetti logic.
                \end{itemize}
            \end{block}
        \end{column}
        \begin{column}{0.48\textwidth}
            \begin{block}{Automated Test Hierarchy}
                \begin{itemize}\footnotesize
                    \item \textbf{Unit Tests (\texttt{tests/unit/}):} Fast, isolated component verification (config loader, pipelines, ONNX quantizer).
                    \item \textbf{Integration Tests (\texttt{tests/integration/}):} Full FastAPI request/response lifecycle and HTMX dashboard flows.
                    \item \textbf{Invariants (\texttt{tests/invariants/}):} Assert mathematical coverage bounds ($1 - \alpha$) and pure Pixi toolchain hygiene.
                    \item \textbf{Coverage Gate:} Strictly enforces $\ge 80\%$ test coverage in CI.
                \end{itemize}
            \end{block}
        \end{column}
    \end{columns}

    \vspace{0.1cm}
    \statusdone

    \script{Our quality assurance quartet runs on every single git commit: Ruff for linting, Mypy for static types, Complexipy for readability, and Pytest enforcing over 80 percent code coverage.}
    \qa{Why enforce Cognitive Complexity below 15?}{Because high cognitive complexity directly correlates with production bugs, unhandled edge cases, and difficult refactorings in mission-critical machine learning systems.}
\end{frame}

% ---------------------------------------------------------------------
% Slide 11: Summary & Architectural Takeaways
% ---------------------------------------------------------------------
\begin{frame}{Chapter Summary}{MLOps Stack \& Storage Architecture Takeaways}
    \begin{block}{Core Storage \& Toolchain Principles}
        \begin{itemize}\small
            \item \textbf{Polyglot Storage Strategy:} Specializing storage tiers---Parquet for metadata, pre-computed tensors for embeddings, SQLite for optimization, and S3 for artifacts---maximizes end-to-end throughput.
            \item \textbf{Zero-Leakage Partitioning:} The Epistemic Diamond splits data into 4 strictly isolated tensors, providing the foundation for distribution-free conformal safety.
            \item \textbf{Three-Tier Object Hierarchy:} Combines cloud public registries, local NVMe edge caches, and immutable DVC remotes to eliminate bandwidth bottlenecks and egress costs.
            \item \textbf{Deterministic Hermeticism:} Pure Pixi environments backed by Rattler caching and strict static gates ensure 100\% reproducible production deployments.
        \end{itemize}
    \end{block}

    \vspace{0.2cm}
    \begin{center}
        \textcolor{deepblue}{\textbf{Next: Chapter 05 --- The Epistemic Diamond DVC DAG \& DagsHub Data Pipeline}}
    \end{center}

    \script{To wrap up Chapter 4: our storage tier transforms what could be a 50-terabyte disk nightmare into a sub-second analytical and training engine, backed by immutable content hashing and rigorous static quality gates.}
    \qa{How does the multi-tier storage architecture support continuous retraining?}{By storing pre-computed foundation embeddings in local SeaweedFS S3, scheduled retraining loops fetch feature vectors in milliseconds and fit new classification heads in under 10 seconds.}
\end{frame}
